
\documentclass{article}
\usepackage{colortbl}
\usepackage{makecell}
\usepackage{multirow}
\usepackage{supertabular}

\begin{document}

\newcounter{utterance}

\centering \large Interaction Transcript for game `wordle\_withclue', experiment `medium\_frequency\_words\_clue\_no\_critic', episode 20 with reka{-}flash{-}3.1{-}t0.0.
\vspace{24pt}

{ \footnotesize  \setcounter{utterance}{1}
\setlength{\tabcolsep}{0pt}
\begin{supertabular}{c@{$\;$}|p{.15\linewidth}@{}p{.15\linewidth}p{.15\linewidth}p{.15\linewidth}p{.15\linewidth}p{.15\linewidth}}
    \# & \multicolumn{2}{c}{Player} && \multicolumn{2}{c}{Game Master} \\
    \hline

    \theutterance \stepcounter{utterance}  
    & & & \multicolumn{4}{p{0.6\linewidth}}{
        \cellcolor[rgb]{0.9,0.9,0.9}{
            \makecell[{{p{\linewidth}}}]{
                \texttt{\tiny{[P1$\langle$GM]}}
                \texttt{You are a language wizard who likes to guess words by using the given rules.} \\
\\ 
\texttt{Welcome to Wordle! You have six attempts to guess the target word, a valid English word of five lowercase letters (a{-}z). Please use the tags "explanation:" and "guess:" to provide a concise explanation for each guess.} \\
\\ 
\texttt{To help you make an informed guess, you will receive a clue for the word, such as} \\
\texttt{clue: snowy white.} \\
\\ 
\texttt{Here is an example guess based on the clue:} \\
\texttt{explanation: In the fairy tail Snow White, the girl is killed because she eats a poisoned apple. And the word apple has 5 letters.} \\
\texttt{guess: apple} \\
\\ 
\texttt{After each guess, your answer will be validated, and you will receive feedback indicating which letters are correct (green), which letters are correct but in the wrong position (yellow), and which letters are incorrect (red). This feedback can be useful in determining which letters to include or exclude in your next guess.} \\
\\ 
\texttt{For example, the feedback for "apple" might be:} \\
\texttt{guess\_feedback: a<yellow> p<yellow> p<green> l<yellow> e<red>} \\
\\ 
\texttt{The explanation should contain details about how both the guess\_feedback and clue work together in tandem to formulate a new guess.} \\
\\ 
\texttt{Do not generate any other text.} \\
\\ 
\texttt{Let's begin with your first guess.} \\
\\ 
\texttt{clue: winter wear} \\
            }
        }
    }
     \\ \\

    \theutterance \stepcounter{utterance}  
    & \multicolumn{4}{p{0.6\linewidth}}{
        \cellcolor[rgb]{0.561,0.737,0.561}{
            \makecell[{{p{\linewidth}}}]{
                \texttt{\tiny{[P1$\rangle$GM]}}
                \texttt{<reasoning>} \\
\texttt{Let's tackle this Wordle with the clue "winter wear". The target is a five{-}letter word. First, I need to think about what "winter wear" refers to. Winter wear typically includes clothing items like coats, jackets, sweaters, scarves, hats, etc. Since it's a five{-}letter word, maybe something shorter like "coat" (4 letters) is too short. "Sweat" is 5 letters. "Scarf" is 5 letters. "Hats" is 4. "Jacket" is 6. "Boots" is 5. "Gloves" is 6. "Mittens" is 7. "Parka" is 5. "Sweat" is 5. "Fleece" is 6. "Socks" is 5. "Cozy" is 5. "Warm" is 4. "Flannel" is 8. "Puffer" is 6. "Down" is 4. "Layer" is 6. "Thick" is 5. "Peacoat" is 7. "Sweater" is 7. "Windbreaker" is 11. "Insulated" is 11. "Winterwear" itself is 10 letters, but the target is 5 letters.} \\
\\ 
\texttt{So possible candidates are "coat", "scarf", "sweat", "boots", "socks", "cozy", "warm", "thick", "parka", "fleece". Let's check which of these are valid 5{-}letter words. "Coat" is valid. "Scarf" is valid. "Sweat" is valid. "Boots" is valid. "Socks" is valid. "Cozy" is valid. "Thick" is valid. "Parka" is valid. "Fleece" is valid. "Warm" is valid. } \\
\\ 
\texttt{Now, considering the clue "winter wear", maybe the word is something that's specifically associated with winter clothing. Let's think of common winter wear items. Coats, scarves, hats, gloves, boots, socks, sweaters, parkas, etc. Since the word is five letters, "scarf" is 5 letters. "Coat" is 4. "Sweat" is 5. "} \\
            }
        }
    }
    & & \\ \\

    \theutterance \stepcounter{utterance}  
    & & & \multicolumn{2}{p{0.3\linewidth}}{
        \cellcolor[rgb]{0.9,0.9,0.9}{
            \makecell[{{p{\linewidth}}}]{
                \texttt{\tiny{[GM$|$GM]}}
                \texttt{Error: The response should always start with the keyword 'explanation:'} \\
            }
        }
    }
    & & \\ \\

    \theutterance \stepcounter{utterance}  
    & & & \multicolumn{2}{p{0.3\linewidth}}{
        \cellcolor[rgb]{0.9,0.9,0.9}{
            \makecell[{{p{\linewidth}}}]{
                \texttt{\tiny{[GM$|$GM]}}
                \texttt{game\_result = ABORT} \\
            }
        }
    }
    & & \\ \\

\end{supertabular}
}

\end{document}
